\section{Introduction}
\label{sec:intro}

\subsection{Motivation}
Online social networks are organised into communities: groups of users who can reach
one another easily and who tend to share views and interests. A false message usually
starts inside one such group. Its reach depends largely on whether it crosses into
other groups, and that crossing happens through comparatively few accounts and links
that connect otherwise separate parts of the network~\citep{shah2024telegram}.
Intervening at those crossings, by removing, throttling or labelling the content that
passes through them, is the premise behind community-aware misinformation
containment~\citep{ghoshal2021community}.

The standard way to choose where to intervene is structural. Influence blocking
methods rank nodes by degree, PageRank, betweenness or bridging scores computed from
the graph alone~\citep{budak2011limiting,xie2025blocking}. Structure, however, only
says who is connected to whom. It does not say who is likely to \emph{accept} a
message. A user whose friends are all in one community but who writes like members of
another is structurally ordinary and invisible to these scores, yet may be exactly the
point where a message from that other community gets in.

\subsection{Idea}
The project builds two separate views of the same users:
\begin{itemize}
  \item a \textbf{topology view}, in which communities are found from the social graph
        alone; and
  \item a \textbf{content view}, in which a classifier that reads only each user's own
        text predicts which of those communities the user belongs to.
\end{itemize}
The content view receives no edges or graph-derived features as input; the topology
communities are used only as its training targets. Where the two views agree, the two
sources of evidence about the user are consistent, although both could still be wrong.
Where they disagree, the user's edges say one thing and their words say another. We call this
\emph{structure--content disagreement} and score it per node.

The idea follows the Ambiguity-Aware Sampling (AAS) framework of
\citet{sani2026aas}, which runs two complementary clustering methods on the same data
and treats the regions where they disagree as the most informative places to spend
an annotation budget. We carry that principle from image embeddings to social
networks. Our two views are structure and content, not two clustering algorithms on
one embedding. The question is whether the disagreement points to useful intervention
targets.

\subsection{Objectives}
\begin{enumerate}
  \item Build a topology view and a content view, where the content view takes no
        edges or graph-derived features as input and uses the topology communities
        only as cross-fitted training targets, and a per-node disagreement score
        between them.
  \item Test whether the score detects nodes whose content and structure genuinely
        conflict, on synthetic networks with planted conflicts and on a real network
        with party labels (TIMME).
  \item Test, by simulating spread with an unknown source, whether targeting nodes
        or links with high disagreement contains spread better than strong
        structural baselines, and map the conditions under which it does.
  \item (Planned extension) Use the disagreement to choose queries in an active
        learning loop, as in AAS.
\end{enumerate}

\subsection{Scope of this report}
This is a mid-term report. It covers the problem, prior work, the methodology as
implemented, the experimental setup, and the status and plan of the work.
Experimental results are deliberately left out and will be reported in the final
report.

\subsection{Organisation}
\Cref{sec:related} reviews related work and explains how it led to the current idea.
\Cref{sec:method} describes the methodology and its limitations. \Cref{sec:setup}
describes the datasets, parameters, baselines and evaluation measures.
\Cref{sec:status} summarises what has been built and what remains.
